\section{Inventory and P\&L Tracking}
\label{sec:pnl}

Every fill on $A$ and every hedge leg on $B$/$C$ is appended to a
single trade log with columns $(t, \mathrm{side}, p, x, m^{\text{fair}}, \mathrm{cf}, \phi, \mathrm{is\_hedge}, q_{\text{after}})$,
where $\mathrm{cf}$ is the signed USD cash-flow and $\phi$ the fee
paid on that leg. All P\&L figures reported in the backtest
(Section~\ref{sec:backtest}) are derived from this log.
\subsection{Sign and fee convention}

Cash-flows are signed from the market-maker's perspective in USD:
positive when cash flows in, negative when it flows out. For a sell
fill on $A$ at price $p$ and size $x$ (EUR), the cash-flow is
$+p\,x\,(1 - \phi_A^{\text{mak}})$; for a buy fill it is
$-p\,x\,(1 + \phi_A^{\text{mak}})$. Hedge legs cross the opposite
side on $B$ or $C$ and pay the corresponding taker fee. Fees are
therefore \emph{already embedded} in $\mathrm{cf}$. We log $\phi$
separately only to isolate the fee line in the decomposition below
--- it is not subtracted a second time.

\subsection{Realized versus mark-to-market}

At any time $t$ we split the total economic P\&L into what is
locked in and what remains at risk:
\begin{equation}
  \mathrm{MtM}_t \;=\; \underbrace{\sum_{i:\,t_i \le t} \mathrm{cf}_i}_{\text{realized}}
                 \;+\; \underbrace{q_t \cdot m^{\text{fair}}_t}_{\text{unrealized}},
  \label{eq:mtm-split}
\end{equation}
where $q_t$ is the open EUR inventory after the last fill and
$m^{\text{fair}}_t$ is the cross-venue fair mid of
Section~\ref{sec:quoter-mid}. Valuing inventory at the fair mid
rather than at $A$'s own noised mid removes the venue-specific
microstructure noise from the valuation: the unrealized line reflects
only the economic value of the EUR position, not the idiosyncratic
distortion of the venue on which it sits.

\subsection{Inception spread versus mid-to-mid revaluation}

Equation~\eqref{eq:mtm-split} separates \emph{what is settled} from
\emph{what is still open}, but it does not say \emph{why} the P\&L
got to its current value. For that we use a second, orthogonal,
decomposition:
\begin{equation}
  \mathrm{MtM}_t \;=\;
  \underbrace{\sum_{i \in \mathcal{M},\,t_i \le t} s_i\,(p_i - m^{\text{fair}}_{t_i})\,|x_i|}_{\text{inception spread } \mathcal{S}_t}
  \;+\;
  \underbrace{\mathrm{Reval}_t}_{\text{mid-to-mid revaluation}}
  \;-\;
  \underbrace{\mathcal{F}_t}_{\text{fees}},
  \label{eq:inception-reval}
\end{equation}
with $\mathcal{M}$ the set of MM fills on $A$ (hedge legs excluded),
$s_i = +1$ for a sell and $-1$ for a buy, and
$\mathcal{F}_t = \sum_{i \le t}|\phi_i|$ the cumulative fees.
$\mathcal{S}_t$ is the spread we captured \emph{at the moment} of
each fill --- it is the pure market-making edge, measured before any
subsequent market move. $\mathrm{Reval}_t$ is the residual, defined
by \eqref{eq:inception-reval}, and isolates the P\&L generated purely
by $m^{\text{fair}}$ drifting after a fill on the open position it
left us with; it is a directional, non-market-making component.
This is precisely the distinction between inception spread and
mid-to-mid revaluation that the brief asks for.

\subsection{Continuous MtM and final inventory}

Evaluating \eqref{eq:mtm-split} only at fill events gives a sparse
view. For the backtest report we reconstruct MtM on the full
simulation grid: the cumulative cash, fees, and inception spread
columns are forward-filled onto every step, and residual inventory
is marked at the concurrent fair mid. This yields a continuous curve
plotted alongside the discrete realized-cash ladder in
Section~\ref{sec:backtest} and makes the mid-to-mid revaluation
component visible between fills, not just at them. At horizon $T$ we
close the run with the triple $(\text{realized}_T,\, q_T \cdot m^{\text{fair}}_T,\, q_T)$:
the first is the locked-in USD P\&L, the second is the final
inventory value in USD --- the mark at which we would liquidate the
remaining EUR --- and the third is the raw EUR position we carry out.
